\documentclass[UTF8,12pt,a4paper]{ctexart}
\usepackage[a4paper,margin=2.15cm]{geometry}
\usepackage{amsmath,amssymb,bm,booktabs,array,longtable,multirow}
\usepackage{enumitem}
\usepackage{fancyhdr}
\usepackage{lastpage}
\linespread{1.3}
\setlength{\parindent}{2em}
\setlength{\parskip}{0.2em}
\setlist[enumerate]{leftmargin=2.5em,itemsep=0.45em,topsep=0.3em}
\setlist[itemize]{leftmargin=2.2em,itemsep=0.25em,topsep=0.2em}
\renewcommand\arraystretch{1.18}
\newcommand{\blank}[1][2.4cm]{\underline{\hspace{#1}}}
\newcommand{\dd}{\mathrm{d}}
\newcommand{\ee}{\mathrm{e}}
\newcommand{\jj}{\mathrm{j}}
\newcommand{\prob}{\mathsf{P}}
\pagestyle{fancy}
\fancyhf{}
\cfoot{第\thepage 页 / 共\pageref{LastPage}页}
\renewcommand{\headrulewidth}{0pt}
\begin{document}

\begin{center}
    {\zihao{2}\bfseries 误差理论与数据处理试卷（二）}\par
    \vspace{0.8em}
\end{center}
\vspace{0.6em}

\section*{一、填空题答案}
\begin{enumerate}[label=\arabic*、]
\item 单峰性、有界性、对称性、抵偿性。
\item $4.510$，$6.379$。
\item $0.3\,\mathrm{N/cm^2}$。
\item $-0.001$，实验对比，$0.001$，$10.003$。
\item $3\sigma$，$10$。
\item $\delta_D=\dfrac{\sqrt{2}\,\delta_V}{\pi Dh}$，$\delta_h=\dfrac{2\sqrt{2}\,\delta_V}{\pi D^2}$。
\item $50\mu\Omega=0.00005\Omega$，B，$8$。
\item $0.067\,\mathrm{mH}$。
\item 绝对，相对。
\end{enumerate}

\section*{二、选择题答案}
\begin{center}
\begin{tabular}{cccccccccc}
\toprule
1&2&3&4&5&6&7&8&9&10\\
\midrule
B&A&A&C&C&B&A&C&C&D\\
\bottomrule
\end{tabular}
\end{center}

\section*{三、判断题答案}
\begin{center}
\begin{tabular}{cccccccccc}
\toprule
1&2&3&4&5&6&7&8&9&10\\
\midrule
对&错&错&错&错&对&对&错&对&错\\
\bottomrule
\end{tabular}
\end{center}

\newpage

\section*{四、计算题答案}
\begin{enumerate}[label=\arabic*、]
\item 加权算术平均值为
\begin{equation*}
    \bar{x}=\frac{\sum p_ix_i}{\sum p_i}=\frac{174.7}{17}=10.2765.
\end{equation*}
单位权标准差为
\begin{equation*}
    \sigma_0=\sqrt{\frac{\sum p_i(x_i-\bar{x})^2}{n-1}}
    =\sqrt{\frac{0.3906}{5}}=0.2795.
\end{equation*}
加权算术平均值的标准差为
\begin{equation*}
    \sigma_{\bar{x}}=\frac{\sigma_0}{\sqrt{\sum p_i}}
    =\frac{0.2795}{\sqrt{17}}=0.0678.
\end{equation*}

\item 电流的最可信赖值为
\begin{equation*}
    I=U/R=16.50/4.26=3.873\,\mathrm{A}.
\end{equation*}
电压、电阻平均值的标准不确定度分别为
\begin{equation*}
    u_U=0.1/\sqrt{4}=0.05\,\mathrm{V},\qquad
    u_R=0.04/\sqrt{4}=0.02\Omega.
\end{equation*}
由电压引起的分量为
\begin{equation*}
    u_1=\left|\frac{\partial I}{\partial U}\right|u_U=\frac{1}{R}u_U=0.012\,\mathrm{A}.
\end{equation*}
由电阻引起的分量为
\begin{equation*}
    u_2=\left|\frac{\partial I}{\partial R}\right|u_R=\frac{U}{R^2}u_R=0.018\,\mathrm{A}.
\end{equation*}
考虑相关系数 $\rho_{UR}=-0.36$，合成标准不确定度为
\begin{equation*}
\begin{aligned}
    u_c&=\sqrt{u_1^2+u_2^2+2\rho_{UR}u_1u_2} \\
       &=\sqrt{0.012^2+0.018^2+2\times(-0.36)\times0.012\times(-0.018)}
        \approx0.025\,\mathrm{A}.
\end{aligned}
\end{equation*}
有效自由度为
\begin{equation*}
    \nu_{\mathrm{eff}}=\frac{u_c^4}{u_1^4/3+u_2^4/3}\approx11.
\end{equation*}
当 $P=99\%$ 时，取 $t_{0.01}(11)=3.11$，所以
\begin{equation*}
    U_{99}=3.11u_c=3.11\times0.025=0.078\,\mathrm{A}.
\end{equation*}

\newpage

\item 由
\begin{equation*}
A=\begin{bmatrix}5&-0.2\\2&-0.5\\1&-1\end{bmatrix},\qquad
L=\begin{bmatrix}0.8\\0.2\\-0.3\end{bmatrix}
\end{equation*}
得
\begin{equation*}
A^TA=\begin{bmatrix}30&-3\\-3&1.29\end{bmatrix},
\qquad (A^TA)^{-1}=\frac{1}{29.7}\begin{bmatrix}1.29&3\\3&30\end{bmatrix}.
\end{equation*}
故
\begin{equation*}
    \hat{X}=(A^TA)^{-1}A^TL=\begin{bmatrix}0.182\\0.455\end{bmatrix}.
\end{equation*}
残差平方和 $\sum v_i^2=0.0051$，因此
\begin{equation*}
    \sigma=\sqrt{\frac{0.0051}{3-2}}=0.071.
\end{equation*}
由 $d_{11}=0.0434,d_{22}=1.01$，得
\begin{equation*}
    \sigma_{X_1}=0.071\sqrt{0.0434}=0.015,
    \qquad
    \sigma_{X_2}=0.071\sqrt{1.01}=0.071.
\end{equation*}

\newpage

\item 设经验公式为 $\hat{y}=b_0+bx$。由数据得
\begin{equation*}
\sum x_i=180,\quad \sum y_i=692,\quad \sum x_i^2=3450,\quad
\sum y_i^2=49470,\quad \sum x_iy_i=13032.
\end{equation*}
因此
\begin{equation*}
\bar{x}=18,\quad \bar{y}=69.2,
\end{equation*}
\begin{equation*}
l_{xx}=\sum x_i^2-\frac{1}{n}\left(\sum x_i\right)^2=210,
\quad
l_{yy}=\sum y_i^2-\frac{1}{n}\left(\sum y_i\right)^2=1583.6,
\end{equation*}
\begin{equation*}
l_{xy}=\sum x_iy_i-\frac{1}{n}\left(\sum x_i\right)\left(\sum y_i\right)=576.0.
\end{equation*}
回归系数为
\begin{equation*}
    b=l_{xy}/l_{xx}=2.74,
    \qquad b_0=\bar{y}-b\bar{x}=19.8.
\end{equation*}
故经验公式为
\begin{equation*}
    \hat{y}=19.8+2.74x.
\end{equation*}
方差分析：
\begin{equation*}
    U=bl_{xy}=1578.24,
    \qquad Q=l_{yy}-U=5.36,
    \qquad S=l_{yy}=1583.6.
\end{equation*}
\begin{center}
\begin{tabular}{ccccc}
\toprule
来源 & 平方和 & 自由度 & 方差 & $F$ \\
\midrule
回归 & $U=1578.24$ & $1$ & -- & -- \\
残余 & $Q=5.36$ & $8$ & $Q/8=0.67$ & $F=1578.24/0.67=2356$ \\
总计 & $S=1583.6$ & $9$ & -- & -- \\
\bottomrule
\end{tabular}
\end{center}
因为
\begin{equation*}
    F=2356>F_{0.01}(1,8)=11.26,
\end{equation*}
所以回归方程在 $0.01$ 水平上高度显著。

\end{enumerate}
\end{document}
